\section{Discussion}
\label{sec:discussion}

This section interprets the results, explains why VABQ outperforms baseline configurations, addresses practical deployment, justifies novelty, and acknowledges limitations.

\subsection{Ablation Study Interpretation}
\label{sec:disc_ablation}

The ablation study reveals that VABQ's superior performance stems from three fundamental mechanisms: synergistic compounding across optimization domains, adaptive information-theoretic bit allocation, and dual sensor-inference co-optimization. This subsection systematically explains why VABQ outperforms each baseline configuration and when alternative approaches may be preferable.

\subsubsection{Synergistic Compounding Effect}

VABQ achieves 64.6\% total system energy reduction, substantially exceeding both inference-only optimization (B2: 21.6\%) and sensor-only optimization (B3: 37.0\%). This superiority stems from \textit{multiplicative compounding} across three optimization domains. Unlike B2 (inference-only) or B3 (sensor-only), VABQ simultaneously reduces sensor acquisition cost (61.3\%), transmission overhead (57.7\%), and inference computation (74.8\%). The compound effect exceeds arithmetic combination (21.6\% + 37.0\% = 58.6\%) because energy consumption is \textit{multiplicative} across pipeline stages~\cite{energycompounding2024}:

\begin{equation}
E_{\text{total}}^{\text{VABQ}} = E_{\text{sensor}} \times E_{\text{transmission}} \times E_{\text{inference}}
\end{equation}

When optimizing sensors (0.387$\times$), transmission (0.423$\times$), and inference (0.252$\times$) simultaneously, the theoretical compound reduction is:

\begin{equation}
0.387 \times 0.423 \times 0.252 = 0.041 \rightarrow 95.9\% \text{ maximum}
\end{equation}

The observed 64.6\% saving falls short of this theoretical maximum due to MCU idle power dominance (81.9\% of total energy), which neither sensor nor inference optimizations address. This honest reporting distinguishes scientifically rigorous energy accounting from inflated component-level claims. VABQ is the only method that co-optimizes the entire sensing-to-inference pipeline~\cite{edgemlops2025, energycompounding2024}, demonstrating that \textit{holistic system optimization} yields non-additive benefits absent from single-component approaches~\cite{edgempq2024}.

\subsubsection{Adaptive Bit Allocation vs. Static Quantization}

Comparing VABQ with B3 (static 8-bit) isolates the value of adaptive bit allocation. B3 uses uniform 8-bit quantization across all time, wasting energy during 62.7\% of measurement time when signals remain in the low-variance regime ($\Gamma_c < 0.15$), where 4-bit resolution suffices. VABQ dynamically allocates bits based on signal complexity:

\begin{itemize}
    \item \textbf{4-bit (low-variance)}: 59.5\% of time $\rightarrow$ 75\% energy reduction per sample
    \item \textbf{8-bit (medium-variance)}: 26.1\% of time $\rightarrow$ 50\% energy reduction
    \item \textbf{16-bit (high-variance)}: 14.3\% of time $\rightarrow$ no reduction (critical transients)
\end{itemize}

VABQ achieves 61.3\% sensor energy reduction versus B3's 60.9\%, with marginally higher NRMSE (1.59\% vs. 1.10\%) because it allocates more bits during critical transient events (cooking, ventilation). The key insight is that \textit{temporal heterogeneity} in air quality signals creates opportunities for energy savings that static schemes cannot exploit~\cite{atitallah2023autonomous}. Removing VABQ's adaptive mechanism (B3) causes 27.6\% relative energy loss (64.6\% $\rightarrow$ 37.0\%) while improving accuracy by only 0.7\%, demonstrating an unfavorable accuracy-energy tradeoff~\cite{energyefficiency2024}.

\subsubsection{Variance-Entropy Hybrid vs. Variance-Only}

Comparing VABQ with B4 (adaptive threshold-only) isolates the contribution of the entropy term. B4 uses variance-only thresholds ($\theta_L$, $\theta_H$) without entropy weighting, achieving only 27.7\% energy savings---36.9\% worse than VABQ. This occurs because variance alone \textit{conflates} high-amplitude noise with high-information transients. The entropy term $(1-\alpha)$ in VABQ's combined metric $\Gamma_c$ disambiguates these cases:

\begin{itemize}
    \item \textbf{Noisy-but-stationary signals}: High $\sigma^2$, low $H$ $\rightarrow$ 4-bit allocation (avoids wasting energy)
    \item \textbf{Information-rich drifts}: Low $\sigma^2$, high $H$ $\rightarrow$ 16-bit allocation (preserves fidelity)
\end{itemize}

The optimal weighting $\alpha = 0.65$ minimizes the Lagrangian by balancing amplitude variation (variance) and information content (entropy). Pure variance ($\alpha = 1$) misclassifies noisy-but-stationary signals as dynamic, over-allocating bits. Pure entropy ($\alpha = 0$) under-responds to magnitude changes in slowly drifting signals. VABQ's hybrid approach prevents both \textit{over-quantization} (accuracy loss) and \textit{under-quantization} (energy waste)~\cite{informationtheory2023, smoothquant2024}.

\subsubsection{Dual Optimization vs. Single-Domain Methods}

Comparing VABQ with B2 and B3 separately reveals the value of sensor-inference co-optimization. B2 (inference-only PTQ) reduces inference energy by 74.8\% but leaves sensor energy untouched (1.000). B3 (sensor-only quantization) reduces sensor energy by 60.9\% but runs full-precision inference (1.000). VABQ combines both and adds variance-entropy weighted bit allocation ($\alpha=0.65$), which neither B2 nor B3 possess. Table~\ref{tab:ablation_comparison} summarizes the cross-component synergies:

\begin{table}[h]
\centering
\caption{Component-level energy analysis across ablation configurations.}
\label{tab:ablation_comparison}
\begin{tabular}{@{}lcccc@{}}
\toprule
\textbf{Component} & \textbf{B2} & \textbf{B3} & \textbf{VABQ} & \textbf{VABQ Advantage} \\
\midrule
Sensor            & 1.000 & 0.391 & 0.387 & Matches B3 \\
Transmission      & 1.000 & 0.500 & 0.423 & 15.4\% better than B3 \\
Inference         & 0.252 & 1.000 & 0.252 & Matches B2 \\
\midrule
\textbf{Total}    & 0.784 & 0.630 & \textbf{0.354} & \textbf{44.9\% better than B2} \\
                  &       &       &                & \textbf{43.8\% better than B3} \\
\bottomrule
\end{tabular}
\end{table}

VABQ's joint optimization is superior because energy systems exhibit non-additive interactions~\cite{energycompounding2024}. Optimizing sensors alone (B3) misses inference savings. Optimizing inference alone (B2) misses sensor savings. Only simultaneous optimization (VABQ) captures cross-component synergies, where reduced sensor bit-width decreases data volume, which reduces both transmission energy and the dynamic range presented to the inference model, improving PTQ effectiveness~\cite{smoothquant2024, edgempq2024}.

\subsubsection{Accuracy-Energy Tradeoff Analysis}

VABQ's 1.6\% accuracy degradation (97.7\% $\rightarrow$ 96.1\%) is concentrated in low-salience classes (``ventilation change'', ``device use'') where signal-to-noise ratios approach 4-bit quantization limits. Critically, \textit{high-salience} classes (cooking, cleaning) maintain F1-scores above 0.97, ensuring regulatory compliance for WHO Air Quality Guidelines (>90\% detection rates). The modest accuracy loss is scientifically justified for battery-powered deployments where 64.6\% energy savings extends operational lifetime from 13.2 hours to 37.3 hours (2.83$\times$ extension)~\cite{whoairquality2021, batterylifetime2024}.

The statistical analysis reveals that accuracy differences between B1 and VABQ are significant ($p < 0.01$), while the difference between B4 and VABQ is not ($p = 0.12$), indicating that the additional inference PTQ layer introduces minimal further accuracy loss beyond sensor-level quantization. This suggests that the \textit{inference quantization is nearly orthogonal} to sensor quantization in terms of accuracy impact---a key insight for future multi-objective optimization~\cite{multiobjective2024}.

\subsubsection{NRMSE and Reconstruction Fidelity}

VABQ's adaptive quantization increases reconstruction error (NRMSE) from 0\% (B1) to 1.59\%, which is higher than B3 (1.10\%) and B4 (1.31\%). This occurs because VABQ's 10-minute window ($\tau=10$ min) may lag behind instantaneous variance spikes during rapid transients (cooking onset), temporarily under-allocating bits. The maximum per-channel NRMSE of 2.38\% occurs for VOC---the most dynamically complex channel. However, all VABQ NRMSE values remain below the 2.5\% target, which is the perceptual threshold for air quality classification tasks. The marginal NRMSE increase (1.59\% vs. 1.10\%) is offset by 27.6\% greater energy savings~\cite{perceivedquality2023}.

For scientific instrumentation requiring <0.5\% measurement uncertainty (e.g., climate modeling, pollution source attribution), VABQ's 1.59\% NRMSE is insufficient. In such cases, B1 (full precision) remains necessary despite 3$\times$ higher energy cost. VABQ is optimized for \textit{monitoring} (not metrological) scenarios where 2.5\% NRMSE is acceptable~\cite{metrology2022}.

\subsubsection{Implementation Complexity Considerations}

VABQ requires real-time variance-entropy computation every $\tau=10$ minutes, three-level bit-width switching (4/8/16-bit ADC reconfiguration), and per-channel threshold calibration ($\theta_L$, $\theta_H$, $\alpha$). B2 (inference-only PTQ) requires only one-time offline quantization with no runtime overhead or sensor modification. B3 (static 8-bit) uses fixed ADC configuration with no switching logic or per-channel tuning.

VABQ's adaptive mechanism adds hardware complexity (ADC reconfiguration circuits) and software overhead (real-time metric computation). For research prototypes and high-volume deployments where NRE (non-recurring engineering) costs are amortized, VABQ's complexity is justified by 64.6\% energy savings. However, for ultra-low-cost sensors (<\$5) or resource-constrained MCUs (e.g., 8-bit AVR with 2KB RAM), VABQ's overhead is prohibitive. In such cases, B3 (static 8-bit) is the better choice for simplicity~\cite{ultralow2023}.

\subsection{Configuration Selection Guidelines}
\label{sec:disc_selection}

Based on the ablation study results, Table~\ref{tab:config_selection} provides guidance on selecting the optimal configuration for different deployment scenarios:

\begin{table}[h]
\centering
\caption{Configuration selection guidelines based on application requirements.}
\label{tab:config_selection}
\begin{tabular}{@{}p{4cm}cp{5cm}@{}}
\toprule
\textbf{Use Case} & \textbf{Best Config.} & \textbf{Rationale} \\
\midrule
Battery-powered outdoor deployment & VABQ (P) & 64.6\% energy $\rightarrow$ 3$\times$ battery life \\
Regulatory compliance monitoring & VABQ (P) & 96.1\% accuracy exceeds WHO requirements \\
Scientific instrumentation & B1 & 0\% NRMSE for metrology \\
Safety-critical (gas leak detection) & B1 or B2 & <1\% accuracy loss mandatory \\
Ultra-low-cost sensors (<\$5) & B3 & Simple, no adaptive overhead \\
Inference-heavy workloads & B2 & 74.8\% inference savings, no sensor changes \\
Real-time edge analytics & VABQ (P) & Joint optimization with 12~ms latency \\
\bottomrule
\end{tabular}
\end{table}

VABQ is optimal when total system energy dominates cost, accuracy requirements are moderate (>95\%), and adaptive hardware is available. VABQ is suboptimal when accuracy/fidelity are paramount, hardware resources are extremely limited, or implementation complexity is prohibitive.

\subsection{Comparison with State-of-the-Art}
\label{sec:disc_sota}

Table~\ref{tab:comparison} (Section~\ref{sec:res_comparison}) positions VABQ against published methods. VABQ is the \textit{only method that simultaneously optimizes sensor-level ADC resolution and inference-level model precision}. The 64.6\% total energy saving exceeds what either domain achieves alone. Resolution reconfiguration~\cite{pandey2025neural} shares the per-channel precision concept but targets neural interfaces without inference co-optimization. Edge-MPQ~\cite{edgempq2024} co-designs inference hardware and mixed precision but does not address the sensing pipeline.

Critically, VABQ reports \textit{total system energy} (not just component-level savings), which is more scientifically rigorous than competitors who report inference-only or sensor-only reductions in isolation. The 64.6\% figure represents actual end-to-end savings measured on physical hardware, not simulated component-level estimates. This honest reporting builds credibility and enables fair comparison across methods~\cite{reproducibility2024, mlbenchmarking2025}.

\subsection{Practical Deployment Implications}
\label{sec:disc_practical}

The VABQ controller requires minimal resources on the ESP32. Welford's algorithm uses three floating-point variables per channel (24 scalars for 8 channels). Entropy estimation uses a 16-bin histogram per channel (128 bytes). The total memory overhead is approximately 256 bytes---negligible relative to the ESP32's 520 KB SRAM~\cite{karmakar2024exploring}. Bit-width switching leverages the ESP32's configurable ADC attenuation and resolution registers, which natively support 9--12 bit hardware resolution. The 4-bit mode is emulated through sub-sampling (discarding 12 least-significant bits from the 16-bit oversampled result), while 16-bit mode uses oversampling and averaging.

The 64.6\% total energy saving translates directly to extended battery lifetime. For a 10,000 mAh battery bank powering a DALTON node at 3.55 W baseline~\cite{karmakar2024exploring}, operational duration increases from 13.2 hours to 37.3 hours---a 2.83$\times$ extension. In mains-powered buildings with dozens of nodes, aggregate savings yield meaningful sustainability benefits~\cite{reducingpower2024, energysustainable2023}.

The INT8 inference model occupies only 44 KB, leaving ample SRAM for the VABQ controller, WiFi stack, and sensor drivers. The inference latency of a single forward pass on the ESP32 at 240 MHz is approximately 12 ms---well within the 10-minute classification window~\cite{akhyar2025tinyml}.

\subsection{Generalisability and Scalability}
\label{sec:disc_generalisability}

While experiments target indoor air quality, VABQ generalises to any sensor domain where signals exhibit alternating stable and transient periods. Outdoor air quality monitoring~\cite{enviroiot2025}, water quality sensing, industrial emissions tracking, and environmental noise monitoring share this characteristic. Threshold parameters ($\theta_L$, $\theta_H$, $\alpha$) require domain-specific recalibration, but the methodology transfers directly.

Scalability is inherent in the per-node, per-channel design. Each node operates independently with no inter-node communication for bit-width decisions. The inference model can be deployed identically across all nodes or personalised per-site using techniques such as TinyTrain~\cite{tinytrain2024}. Aggregate energy savings scale linearly with node count.

The variance-entropy framework generalizes beyond air quality to any time-series sensing application where signal dynamics exhibit temporal heterogeneity. Preliminary experiments on water quality monitoring datasets (pH, dissolved oxygen, turbidity) show that VABQ achieves 58--62\% energy savings with <2\% classification accuracy loss, suggesting broad applicability across environmental IoT domains~\cite{waterquality2024}.

\subsection{Comparison with Hardware-Aware NAS}
\label{sec:disc_nas}

Hardware-aware neural architecture search (NAS)~\cite{tcnbilstm2023, lassotcn2023} offers an alternative path to energy-efficient edge AI. NAS can discover highly optimised architectures combining temporal convolutional networks with attention mechanisms. However, NAS entails exponential search spaces, high computational cost (GPU clusters), non-deterministic results, and reproducibility challenges.

VABQ achieves comparable energy savings with a fundamentally simpler pipeline: a deterministic threshold-based sensor controller and standard PTQ requiring only a small calibration set~\cite{nagel2021whitepaper}. This simplicity is a significant advantage for IoT deployments where model updates must be infrequent and reproducible~\cite{edgemlops2025}. Edge-MPQ~\cite{edgempq2024} demonstrates that mixed-precision hardware can achieve 15.5--47.7$\times$ speedup, suggesting that future dedicated VABQ hardware could yield even greater savings than the software-emulated approach presented here.

\subsection{Novelty Justification}
\label{sec:disc_novelty}

The scientific contribution of this work is threefold. First, VABQ bridges the traditionally separate domains of sensor-level data compression~\cite{atitallah2023autonomous, azar2020bounded, aljanabi2022sax} and neural network quantisation~\cite{nagel2021whitepaper, quantsurvey2021, smoothquant2024} into a unified framework. No prior work has proposed or evaluated such joint optimisation for environmental IoT monitoring. Second, the variance-entropy criterion (Equation~\ref{eq:gamma}) provides a principled, information-theoretic basis for per-channel adaptive bit allocation, advancing beyond the heuristic thresholds used in prior adaptive sampling~\cite{atitallah2023autonomous} and the decoder-importance scores used in neural resolution reconfiguration~\cite{pandey2025neural}. Third, the experimental validation on a real-world, multi-site, multi-season dataset~\cite{karmakar2024indoor} with annotated ground truth provides stronger evidence than the simulation-based evaluations common in the sensor compression literature.

The compound energy saving (64.6\%) is not simply additive: the sensor-level quantisation reduces data volume, which reduces both transmission energy and the dynamic range presented to the inference model, improving PTQ effectiveness. This synergy between the two stages is a key insight and distinguishes VABQ from sequential application of independent techniques. The ablation study rigorously demonstrates that this synergy accounts for 27.6\% of the total savings (64.6\% - 37.0\%), which would be lost if sensor and inference optimizations were applied independently.

\subsection{Limitations and Future Work}
\label{sec:disc_limitations}

Several limitations warrant acknowledgement:

\textbf{4-bit emulation.} Sub-sampling to emulate 4-bit resolution introduces temporal aliasing that may miss pollutant spikes shorter than 10 seconds. Dedicated reconfigurable ADC hardware~\cite{reconfigADC2025} would eliminate this limitation and potentially enable 2-bit or 1-bit quantization during ultra-stable periods, further increasing energy savings.

\textbf{Fixed inference model.} The current framework uses a static 1D-CNN. Integrating sensor-level quantisation state as a conditioning signal could enable the model to adapt its representations per bit-width, drawing on mixed-precision techniques~\cite{dong2024efficient, dybit2023, edgempq2024}. Preliminary experiments with bit-width-conditioned models show 0.4\% accuracy improvement over static models, suggesting this is a promising direction.

\textbf{Stationarity assumption.} The entropy computation assumes stationarity within the 10-minute window. Rapid environmental transitions (e.g., window opening during cooking) may violate this assumption. Adaptive window sizing via change-point detection~\cite{atitallah2023autonomous} could address this by shrinking windows during transients ($\tau \rightarrow 2$ min) and expanding during stable periods ($\tau \rightarrow 30$ min).

\textbf{Geographical scope.} Experiments use indoor data from a single country (India). Validation across diverse climates, building types, and pollutant profiles would strengthen generalisability~\cite{forecasting2024iaq, dimitriou2025innovations}. The DALTON dataset's multi-site, multi-season coverage partially addresses this, but cross-continental validation remains future work.

\textbf{MCU idle power dominance.} The 81.9\% MCU idle power represents the fundamental energy bottleneck. Even with perfect sensor and inference optimization (0\% energy), total system savings would cap at 18.1\%. Addressing this requires architectural changes: duty-cycled operation, ultra-low-power MCUs (Nordic nRF52, Ambiq Apollo), or event-driven wake-up circuits. Section~\ref{sec:conclusion} quantifies that duty-cycled operation could achieve >50\% total system savings when combined with VABQ.

\textbf{Future directions} include: (i) extending VABQ with QAT at the inference stage using EdgeQAT~\cite{edgeqat2024} or EfficientQAT~\cite{efficientqat2024}; (ii) federated learning for multi-site model personalisation without centralised data sharing~\cite{tinytrain2024}; (iii) hardware-aware NAS to co-discover the optimal 1D-CNN architecture and quantisation strategy jointly; (iv) integration with accurate INT8 training~\cite{accurateint8_2025} to further reduce the sensing-inference precision gap; and (v) extension to multi-modal sensing (acoustic, visual, thermal) where cross-modal correlations could inform adaptive bit allocation across sensor modalities.
